\infographicchapter{Cardiac MRI (CMR)}{cardiac_mri.png}
\label{ch:infographic-cardiac-mri}

\section{Overview}

Cardiac magnetic resonance imaging (cardiac MRI, or CMR) is a painless, noninvasive heart scan that uses a strong magnet, harmless radio waves, and a computer to create crystal-clear still and moving pictures of your beating heart and major blood vessels. Unlike X-rays or heart CT scans, cardiac MRI uses zero ionizing radiation, making it safe when children or adults need repeat scans over many years. To time each picture to your heartbeat, small sticky ECG patches are placed on your chest, and for many exams a small IV line delivers a gadolinium contrast dye that helps highlight damaged or inflamed heart muscle.

Cardiac MRI is considered the gold standard for measuring heart chamber sizes, pumping strength (ejection fraction), and the exact health of the heart muscle tissue itself. Doctors rely on it to look for permanent heart attack scar versus living, recoverable muscle (called viability using late gadolinium enhancement), pinpoint heart muscle inflammation (myocarditis), diagnose inherited or acquired heart muscle diseases (cardiomyopathies), map complex birth defects (congenital heart disease), and measure blood flow across leaky or narrowed heart valves.

Because the scanner captures your heart in crisp, multi-angle ``movies'' without needing bones or lungs to be out of the way, it often answers questions that an ultrasound (echocardiogram) leaves unclear. At the same time, a cardiac MRI takes longer (typically 30--60 minutes), asks you to lie still inside a noisy tunnel and hold your breath briefly for sets of pictures, and requires careful screening for metal implants, pacemakers, and severe kidney disease before gadolinium dye is given.

\section{Underlying Physical Principles}
\begin{itemize}
\item \textbf{Magnets and Hydrogen Atoms:} A powerful superconducting magnet (usually 1.5 or 3.0 Tesla) gently aligns tiny hydrogen protons found naturally inside the water and fat of your heart muscle, blood, and chest tissues.
\item \textbf{Radio Waves and Natural Tissue Contrast:} Brief pulses of harmless radio waves nudge those hydrogen protons out of alignment. As they relax back into place, healthy heart muscle, fluid swelling (edema), fat, and flowing blood send back distinct radio echoes without any X-ray radiation.
\item \textbf{Heartbeat Synchronization (ECG Gating):} Because the heart beats continuously, sticky chest ECG electrodes synchronize image capture to each cardiac cycle, freezing motion so the computer can assemble smooth, frame-by-frame ``cine'' movies of every beat.
\item \textbf{Late Gadolinium Enhancement (LGE) for Scar and Viability:} An IV dose of gadolinium dye washes quickly out of tight, healthy heart cells over 10--15 minutes, but lingers longer in expanded spaces where cells are swollen, inflamed, or replaced by scar tissue---making scars and fibrosis glow brightly against dark healthy muscle.
\item \textbf{Phase-Contrast Blood Flow and Tissue Mapping:} Specialized magnetic pulses measure the exact speed and direction of moving blood across heart valves and great vessels, while T1, T2, and T2* mapping quantify hidden swelling, diffuse scarring, amyloid protein, or iron overload.
\end{itemize}

\section{Imaging Mechanism}
\begin{enumerate}
\item \textbf{ECG Setup, Coil Placement, and IV Line:} You lie flat on your back on a padded table while a technologist places MRI-safe ECG stickers on your chest, rests a lightweight cushioned receiver coil over your chest, and places a small arm IV if contrast dye or stress medicine is planned.
\item \textbf{Cine Movies and Baseline Tissue Scans:} After quick scout views locate your heart, the scanner records beating-heart cine movies across standard long- and short-axis planes, along with fluid- and fat-sensitive baseline sequences during short 10--15 second breath-holds.
\item \textbf{Blood Flow and Optional Stress Perfusion:} Phase-contrast sequences measure blood flow volume and velocity through the heart valves and aorta; if checking for coronary artery blockages, a short-acting stress medicine widens heart arteries while first-pass contrast dye tracks blood delivery.
\item \textbf{Delayed Scar Imaging (Late Gadolinium Enhancement):} Roughly 10--15 minutes after the gadolinium injection, specialized inversion-recovery sequences darken healthy heart muscle so any heart attack scar, myocarditis, or cardiomyopathy fibrosis shines clearly.
\item \textbf{Computer Reconstruction and Cardiologist/Radiologist Review:} Advanced software calculates chamber volumes, ejection fraction, valve leakage percentages, and tissue maps, which a specialized cardiac radiologist or cardiologist interprets to guide your care.
\end{enumerate}

\section{Key Technical Parameters}
\noindent
\begin{tabularx}{\textwidth}{@{}l X@{}}
\toprule
\textbf{Parameter} & \textbf{Typical Value / Description} \\
\midrule
Physical Principle & Strong magnetic field, radiofrequency pulses, ECG heartbeat gating, and timed gadolinium washout or blood-flow phase tracking \\
Energy / Signal Source & Superconducting magnet, gradient coils, radiofrequency pulses, chest phased-array receiver coil, and ECG gating leads \\
Typical Frequency / Energy & 1.5 Tesla or 3.0 Tesla clinical magnetic field strength (about 64 MHz or 128 MHz hydrogen frequency) \\
Spatial Resolution & Sharp soft-tissue detail (about 1.2--1.8 mm in-plane) with 6--8 mm cine and scar slices (1--2 mm isotropic for 3D vessels) \\
Typical Scan Duration & 30--60 minutes for most exams; up to 75--90 minutes for complex congenital heart disease or stress perfusion studies \\
Contrast Agent & IV modern macrocyclic gadolinium dye (ACR Group II) for scar, myocarditis, and perfusion; none needed for pure valve flow, anatomy, or iron scans \\
Ionizing Radiation & None (uses safe magnetic fields and radio waves) \\
\bottomrule
\end{tabularx}

\section{Common Symptoms That Prompt This Scan}
\begin{itemize}
\item \textbf{Unexplained Chest Pain or Suspected Myocarditis:} Chest discomfort, palpitations, or elevated heart blood markers (troponin) after a viral illness or when coronary arteries look clear, raising concern for heart muscle inflammation.
\item \textbf{Shortness of Breath, Swelling, or Heart Failure Symptoms:} Reduced pumping strength on an echocardiogram, unexplained leg swelling, or fatigue requiring precise evaluation of heart muscle weakness (cardiomyopathy) or prior silent heart attack damage.
\item \textbf{Fainting, Palpitations, or Dangerous Heart Rhythms:} Unexplained fainting spells or abnormal heart rhythms (arrhythmias), especially in athletes or people with a family history of sudden cardiac arrest or inherited heart muscle disease.
\item \textbf{Heart Murmurs or Valve Leakage:} Follow-up of an inconclusive echocardiogram to accurately measure how much blood is leaking backward across an aortic or mitral valve, or to check congenital heart repairs.
\end{itemize}

\section{Clinical Usages}

\subsection{Primary Usages}
\begin{itemize}
\item \textbf{Heart Muscle Scar and Viability (LGE):} After a heart attack or in coronary artery disease, late gadolinium enhancement distinguishes thin, permanently scarred muscle from stunned but living muscle, showing whether restoring blood flow with a stent or bypass surgery is likely to improve pumping strength.
\item \textbf{Myocarditis and Pericardial Disease:} Detects active heart muscle swelling, tiny blood vessel leakiness, and characteristic outer-wall or mid-wall scar patterns (Lake Louise criteria) in suspected myocarditis, and clearly shows thickened or inflamed heart-sac tissue in constrictive pericarditis.
\item \textbf{Cardiomyopathies and Infiltrative Diseases:} Diagnoses and risk-stratifies hypertrophic, dilated, and arrhythmogenic right ventricular cardiomyopathies, and spots distinctive patterns of cardiac amyloidosis, sarcoidosis, or iron overload (using T2* mapping in thalassemia or frequent blood transfusions).
\item \textbf{Congenital Heart Disease:} Maps complex heart chambers, surgical baffles, shunts, and great vessels in children and adults born with heart defects (such as repaired tetralogy of Fallot or coarctation of the aorta) with zero radiation over lifelong follow-up.
\item \textbf{Heart Valve Blood Flow and Aortic Disease:} Measures exact regurgitant blood volume and fraction in leaky valves (especially aortic and pulmonary valves), assesses valve narrowing, and tracks aneurysms or dissections of the thoracic aorta.
\end{itemize}

\subsection{Secondary Usages}
\begin{itemize}
\item \textbf{Stress Perfusion for Coronary Artery Disease:} Uses a safe vasodilator medicine (such as regadenoson or adenosine) and first-pass contrast dye to spot regions of heart muscle receiving reduced blood flow during stress, without X-rays or nuclear tracers.
\item \textbf{Heart Masses and Blood Clots:} Distinguishes harmless fat, cysts, or blood clots (thrombi) inside heart chambers from true cardiac tumors such as atrial myxoma or sarcoma.
\item \textbf{Pre-Procedure Planning and Cardio-Oncology:} Maps pulmonary veins and scar pathways before arrhythmia ablation or device implantation, and monitors subtle heart muscle changes during cancer chemotherapy.
\end{itemize}

\section{Advantages}
\begin{itemize}
\item \textbf{Zero Ionizing Radiation:} Completely free of X-rays, making it ideal for children, young adults, and patients with congenital heart disease or cardiomyopathies who need repeat scans over their lifetime.
\item \textbf{Unmatched Heart Muscle Tissue Detail:} Uniquely visualizes heart attack scar, active inflammation, diffuse fibrosis, fat, amyloid protein, and iron inside living heart muscle without needing a biopsy.
\item \textbf{Gold Standard for Chamber Size and Pumping Function:} Measures right and left ventricular volumes, wall thickness, and ejection fraction more accurately and reproducibly than ultrasound, unaffected by body build or lung air.
\item \textbf{Exact Blood Flow Quantification:} Phase-contrast imaging directly measures forward and backward blood flow across heart valves and shunts, helping time valve repair or replacement accurately.
\item \textbf{All-in-One Heart Assessment:} Evaluates heart anatomy, muscle motion, valve flow, great vessels, pericardium, and tissue health in a single noninvasive visit.
\end{itemize}

\section{Disadvantages}
\begin{itemize}
\item \textbf{Longer Exam Time and Breath-Holding:} Takes 30--60 minutes (longer than CT or echocardiography) and usually requires repeated 10--15 second breath-holds, which can feel tiring for very breathless patients.
\item \textbf{Irregular Heartbeat Blur:} Very fast or erratic heart rhythms (such as uncontrolled atrial fibrillation) can make standard ECG-gated images blurry unless specialized single-beat sequences are used.
\item \textbf{Metal and Pacemaker Restrictions:} Non-MRI-safe devices or metal fragments are prohibited, and even MRI-conditional pacemakers require specialized cardiology coordination and can still cast small dark shadows over nearby heart walls.
\item \textbf{Gadolinium Dye Precautions:} Scar, inflammation, and stress perfusion studies rely on IV gadolinium contrast, which is generally avoided in pregnancy and weighed carefully in severe kidney failure.
\item \textbf{Noise, Claustrophobia, and Higher Cost:} Requires lying flat inside a narrow, loud scanner tunnel, costs more, and is less widely available outside specialized hospital imaging centers.
\end{itemize}

\section{Contraindications}

\subsection{Absolute}
\begin{itemize}
\item \textbf{Non-MRI-Safe Implants and Metal Fragments:} Older non-MRI-conditional pacemakers or defibrillators (unless scanned at a specialized center under a strict safety protocol), magnetic brain aneurysm clips, cochlear implants, metallic eye fragments, and certain older mechanical heart valves (such as early Starr--Edwards caged-ball valves) cannot enter the magnet room.
\end{itemize}

\subsection{Relative}
\begin{itemize}
\item \textbf{MRI-Conditional Heart Devices:} Modern MRI-conditional pacemakers, ICDs, and loop recorders can usually be scanned safely after a device check and temporary programming into MRI-safe mode, though wires may slightly blur nearby images.
\item \textbf{Severe Kidney Disease or Past Dye Allergy:} In severe kidney failure or dialysis, doctors carefully weigh risks before giving gadolinium dye (preferring modern macrocyclic agents or skipping contrast when anatomy or valve flow alone suffices). A prior severe gadolinium allergy may require premedication or a non-contrast protocol.
\item \textbf{Pregnancy, Severe Claustrophobia, or Unstable Condition:} Contrast dye is generally avoided during pregnancy unless essential (non-contrast CMR can still be performed safely). Severe anxiety, inability to lie flat, or unstable medical status may limit tolerance or require mild sedation or an alternative test.
\end{itemize}

\section{Patient Safety Guidelines}
\begin{itemize}
\item \textbf{Strict Magnet Room Screening:} Every patient completes written and verbal screening for pacemakers, defibrillators, artificial heart valves, stents, clips, insulin pumps, medication patches, and past metal-working eye injuries before entering the scanner zone.
\item \textbf{Specialized Pacemaker and Device Protocols:} Patients with MRI-conditional pacemakers or defibrillators are checked by a cardiology device team before and after the scan, with heart rhythm and oxygen continuously monitored during imaging.
\item \textbf{Safe Gadolinium Contrast Use:} Centers use stable ACR Group II macrocyclic gadolinium agents to keep kidney and tissue risks exceedingly low, avoid contrast during pregnancy when possible, and follow ACR guidance confirming breastfeeding mothers do not need to interrupt nursing.
\item \textbf{Hearing Protection, Skin Care, and Call Button:} Earplugs or headphones protect against loud knocking sounds, foam padding prevents skin-to-skin heating loops, and you hold a squeeze-ball call button the entire scan.
\item \textbf{Stress Medication Monitoring:} If stress perfusion medicine is used, nurses and doctors continuously track your blood pressure, ECG heart tracing, and symptoms, with reversal medicine immediately on hand.
\end{itemize}

\section{Prerequisite Tests}
\begin{itemize}
\item \textbf{Prior Echocardiogram and ECG:} Most patients first have a resting ECG and heart ultrasound (echocardiogram) so the cardiac MRI protocol can be tailored to the exact clinical question.
\item \textbf{Kidney Function Blood Test When Indicated:} While modern macrocyclic gadolinium agents rarely require routine labs in healthy patients, a recent creatinine and eGFR blood test is typically checked if you have known kidney disease, are over 60, or per hospital policy.
\item \textbf{Implant Safety Records and Eye X-Ray if Needed:} Bring manufacturer cards for any pacemaker, artificial heart valve, or stent; anyone with past welding or sheet-metal eye exposure may need a quick eye X-ray first.
\item \textbf{Medication Review Before Stress Scans:} If your exam includes stress perfusion imaging, your care team reviews heart medicines (such as beta-blockers or aminophylline-like drugs) and confirms caffeine avoidance.
\end{itemize}

\section{Preparing for the Scan}
\begin{itemize}
\item \textbf{Caffeine and Food Rules:} For routine cardiac MRI without stress medicine, you can usually eat a light meal and take regular medications unless told to fast for 4--6 hours. If a stress perfusion cardiac MRI is scheduled, avoid all caffeine (coffee, tea, soda, chocolate, and energy drinks) for 12--24 hours beforehand, as caffeine blocks the stress medicine.
\item \textbf{Remove All Metal and Change Into a Gown:} Leave jewelry, watches, hairpins, eyeglasses, belts, credit cards, hearing aids, and foil-backed medicine patches outside the scan room, and change into metal-free clothing or a hospital gown.
\item \textbf{Bring Implant Cards and Prior Heart Studies:} Carry device identification cards for any pacemakers, stents, or artificial valves, and bring outside echocardiogram, CT, or angiogram reports if scanned at a new center.
\item \textbf{Plan for Claustrophobia or Breath-Holds:} Let your doctor know ahead of time if enclosed spaces bother you or if lying flat causes shortness of breath; mild calming medicine (with a ride home arranged) or wide-bore scanners can help.
\end{itemize}

\section{During Scan}
\begin{itemize}
\item \textbf{Positioning, ECG Leads, and Chest Coil:} You lie on your back on a cushioned table while ECG stickers are attached to your chest (sometimes a small patch of chest hair is trimmed so leads stick cleanly) and a light receiver blanket or coil rests gently across your chest.
\item \textbf{Following Short Breath-Hold Prompts:} A prerecorded voice over the intercom will ask you to take a gentle breath in, breathe out, and hold your breath for about 10--15 seconds at a time while images are captured, followed by calm breathing breaks.
\item \textbf{Loud Knocking Noises and Contrast Feel:} Wearing earplugs or headphones, you will hear loud rhythmic tapping, buzzing, and thumping sounds as the magnet switches pulses. When gadolinium enters your IV, you may feel brief coolness in your arm or a mild metallic taste; during stress perfusion, brief flushing or warmth fades within a minute or two.
\item \textbf{Continuous Communication and Typical Length:} You hold a squeeze-ball call button the whole time and can speak with the technologist between sequences. Most scans finish in 30--60 minutes.
\end{itemize}

\section{Normal Appearance}
\begin{itemize}
\item \textbf{Smooth, Evenly Thick Muscle Walls:} On cine movies, the left and right heart chambers squeeze strongly and symmetrically, with normal muscle thickness, smooth valve opening and closing, and normal ejection fraction.
\item \textbf{Uniformly Dark Healthy Muscle on Scar Scans:} After late gadolinium enhancement, normal healthy heart muscle looks completely dark and uniform, with no bright white scar patches in any wall layer.
\item \textbf{Even First-Pass Blood Flow and Clear Pericardium:} If stress perfusion is done, contrast dye washes rapidly and evenly throughout all heart muscle regions, and the outer heart sac (pericardium) appears as a razor-thin dark line without trapped fluid.
\item \textbf{Smooth Valve Flow and Normal Great Vessels:} Blood flows smoothly forward across each heart valve with minimal or no backward leak, and the aorta and pulmonary arteries show normal caliber and smooth walls.
\end{itemize}

\section{Observing for Abnormal Patterns}
\begin{itemize}
\item \textbf{Subendocardial to Full-Thickness Scar (Heart Attack):} Bright white late gadolinium enhancement starting along the inner lining (subendocardium) and matching a coronary artery territory marks heart attack scar; scar covering less than 50\% of wall thickness suggests viable muscle that may recover after restoring blood flow.
\item \textbf{Non-Ischemic Mid-Wall or Outer-Wall Enhancement:} Bright patches in the middle or outer (epicardial) layers---especially along the lower-side wall---point toward myocarditis, while mid-wall stripes or insertion-point spots appear in dilated or hypertrophic cardiomyopathies.
\item \textbf{Diffuse Swelling, Amyloid, or Iron Overload:} High T2 signals reveal acute water swelling (edema) in myocarditis or recent heart attacks, global difficulty darkening healthy muscle with high T1/extracellular volume flags cardiac amyloidosis, and abnormally short T2* times reveal excess iron.
\item \textbf{Valve Leaks, Shunts, and Congenital Connections:} Dark turbulent flow jets on cine movies paired with elevated backward flow percentages on phase-contrast images quantify leaky valves, while 3D views map septal holes, abnormal vein connections, or aortic narrowing.
\item \textbf{Pericardial Thickening or Heart Masses:} Bright, thickened outer sac tissue with abnormal septal bounce indicates constrictive pericarditis, while non-enhancing dark spots in a weak chamber apex point to blood clots rather than vascular tumors.
\end{itemize}

\section{Dealing with Uncertainties}
\begin{itemize}
\item \textbf{Irregular Heartbeats and Breathing Motion:} Frequent skipped beats, atrial fibrillation, or incomplete breath-holds can blur cine movies or scar images; technologists counter this using real-time single-beat imaging, motion correction, or free-breathing sequences.
\item \textbf{Implant and Wire Shadows:} Pacemaker batteries, sternal wires, or mechanical valves create local dark magnetic shadows; specialized wideband scar sequences help clear away device artifacts so surrounding heart muscle stays readable.
\item \textbf{Borderline Scar Timing or Subtle Swelling:} Selecting the wrong inversion delay can make healthy muscle look gray instead of black; technologists run quick calibration checks (TI scouts) or repeat views with phase-sensitive reconstruction to avoid false scar calls.
\item \textbf{Correlating With Other Heart Tests:} When a tissue pattern overlaps between conditions (such as healed myocarditis versus non-ischemic cardiomyopathy), cardiologists combine CMR findings with blood biomarkers, ECGs, echocardiograms, coronary imaging, or genetic testing.
\end{itemize}

\section{Augmenting Methods}
\begin{itemize}
\item \textbf{Parametric T1, T2, and Extracellular Volume (ECV) Mapping:} Quantitative pixel-by-pixel color maps uncover diffuse scarring, subtle inflammation, or early amyloid infiltration even when standard scar images look completely normal.
\item \textbf{4D Flow CMR:} Captures three-dimensional blood flow direction and velocity over the entire heartbeat in one volume scan, ideal for complex congenital heart defects, multi-valve leaks, and aortic flow swirls.
\item \textbf{Pharmacologic Stress Perfusion and Feature Tracking:} Combines vasodilator stress testing with myocardial strain analysis (tracking tiny muscle fiber shortening) to detect early heart trouble before pumping fraction drops.
\item \textbf{Coronary CT Angiography and Nuclear PET:} Coronary CT maps narrowed heart arteries and calcium directly, while cardiac FDG-PET adds metabolic inflammation tracking in tricky cases of cardiac sarcoidosis or endocarditis.
\end{itemize}

\section{Comparison with Alternative Modalities}
\begin{itemize}
\item \textbf{Versus Echocardiography (Heart Ultrasound):} Ultrasound is fast, portable, bedside-friendly, inexpensive, and superior for watching fast-moving valve leaflets in real time. Cardiac MRI wins for exact right- and left-ventricle volume measurement, crystal-clear views in patients with poor ultrasound windows, and unique scar/inflammation tissue characterization.
\item \textbf{Versus Cardiac CT Angiography:} Cardiac CT takes only minutes and excels at directly picturing coronary artery plaques, calcium, and stents. Cardiac MRI avoids X-ray radiation and iodine dye, provides dynamic beating-heart tissue characterization (LGE and mapping), and measures valve flow directly.
\item \textbf{Versus Nuclear Stress Scans (SPECT / PET):} Nuclear scans are widely available and helpful in patients who cannot enter an MRI magnet. Cardiac MRI delivers higher spatial detail with zero ionizing radiation and simultaneously checks chamber function, valves, and non-ischemic causes like myocarditis.
\end{itemize}

\section{Follow-up Steps}
\begin{itemize}
\item \textbf{Immediate Post-Scan Care:} Your ECG stickers and IV line are removed right away, and you can immediately return to normal meals, driving, and daily activities (unless you took a relaxing sedative, which requires someone else to drive you home). Breastfeeding mothers may resume nursing immediately with no interruption.
\item \textbf{Quantitative Report Review:} A board-certified cardiac radiologist or cardiologist reviews the cine movies, valve flow numbers, and scar maps, sending a detailed report with chamber volumes, ejection fraction, and tissue findings to your ordering doctor.
\item \textbf{Tailored Treatment Decisions:} Depending on results, your care team may adjust heart failure or anti-inflammatory medicines, clear an athlete to safely return after myocarditis recovery, decide whether a coronary stent, bypass, valve repair, or implanted defibrillator is needed, or arrange family genetic screening.
\item \textbf{Long-Term Radiation-Free Monitoring:} When ongoing checkups are needed---such as tracking congenital heart repairs, aortic aneurysms, or iron-overload therapy---repeat cardiac MRIs safely monitor changes over time without adding lifetime radiation exposure.
\end{itemize}

\section{Key References}
\begin{itemize}
\item American College of Radiology (ACR), North American Society for Cardiovascular Imaging (NASCI), and Society for Pediatric Radiology (SPR). \textit{ACR--NASCI--SPR Practice Parameter for the Performance and Interpretation of Cardiac Magnetic Resonance Imaging (MRI)}. Reston, VA: ACR; revised 2023.
\item American College of Radiology (ACR). \textit{ACR Manual on MR Safety} and \textit{ACR Manual on Contrast Media}. Reston, VA: ACR; 2024.
\item Schulz-Menger J, Bluemke DA, Bremerich J, et al. Standardized image interpretation and post-processing in cardiovascular magnetic resonance --- 2020 update: Society for Cardiovascular Magnetic Resonance (SCMR) Board of Trustees Task Force on Standardized Post-Processing. \textit{Journal of Cardiovascular Magnetic Resonance}. 2020;22(1):19.
\item U.S. National Library of Medicine, MedlinePlus. \textit{Heart MRI}. MedlinePlus Medical Encyclopedia. Available at: \url{https://medlineplus.gov/ency/article/003795.htm}.
\item National Heart, Lung, and Blood Institute (NHLBI), National Institutes of Health. \textit{Heart Tests}. Available at: \url{https://www.nhlbi.nih.gov/health/heart-tests}.
\item Stanford Health Care. \textit{Cardiac MRI}. Available at: \url{https://stanfordhealthcare.org/medical-tests/c/cardiac-mri.html}.
\item Radiopaedia.org. \textit{Cardiac MRI}. Radiopaedia. Available at: \url{https://radiopaedia.org/articles/cardiac-mri}.
\end{itemize}

\section{Further Reading}
\begin{itemize}
  \item Wikipedia contributors. \textit{Cardiac magnetic resonance imaging}. Wikipedia, The Free Encyclopedia. Available at: \url{https://en.wikipedia.org/wiki/Cardiac_magnetic_resonance_imaging}.
\end{itemize}
